\section{Conclusion}
\label{sec:conclusion}

SP500Vol-Bench links 144{,}129 survivorship-free, point-in-time S\&P~500 filings to forward realised volatility, and its incremental-value protocol credits text only for the out-of-sample QLIKE it removes beyond a reference interval of recalibrated price forecasts, under day-clustered Diebold--Mariano inference, Holm correction, and placebo gates. The protocol's first output is a rigorous near-null: no text or fusion model beats HAR-RV out of sample (0 of 180 standalone comparisons), and the 38 of 69 combination cells genuine on the seed-ensemble primary against the most generous single-HAR reference fall to 0 of 69 through the reference interval---a zero-text firm-identity forecast reproduces the gain, the five-model maximal price pool absorbs the remainder, and the two survivor sets are disjoint. The residual---$+$0.2--0.5\% QLIKE in a handful of event-driven 8-K cells---is cross-sectional rather than regime-timing, economically nil, and does not replicate across LLM families. The artefact that outlasts this verdict is the benchmark and the reusable protocol: any claimed disclosure-text gain can now be required to survive the same controls. Our near-null is scoped: it concerns large-cap U.S.\ equities (S\&P~500, 2010--2025), daily-return realised volatility at 5-, 10- and 20-day horizons, against strong price baselines; small-cap or less-efficient markets, intraday or option-implied targets, and other document classes may behave differently. Whether disclosure text clears the same bar there is the question SP500Vol-Bench is built to ask.
